\documentclass[12pt]{article}
\usepackage[utf8]{inputenc}
\usepackage[spanish]{babel}
\usepackage{geometry}
\usepackage{enumitem}
\usepackage{fancyhdr}

\geometry{a4paper, margin=1in}

\title{Examen de Física Básica: Cinemática y Newton - Forma A}
\author{Nombre: \underline{\hspace{8cm}} \quad Matrícula: \underline{\hspace{3cm}}}
\date{Fecha: \underline{\hspace{4cm}}}

\pagestyle{fancy}
\fancyhf{}
\rhead{Examen Forma A}
\lhead{Física Básica}
\cfoot{\thepage}

\begin{document}

\maketitle

\section*{I. SECCIÓN TEÓRICA (4 Puntos)}

\begin{enumerate}[label=\arabic*.]
    \item \textbf{La rama de la física que se encarga de la descripción del movimiento, sin considerar las causas que lo producen, es:}
    \begin{enumerate}[label=\alph*)]
        \item Dinámica
        \item Estática
        \item Cinemática
        \item Mecánica Cuántica
    \end{enumerate}

    \item \textbf{La línea imaginaria que describe un objeto durante su movimiento se conoce como:}
    \begin{enumerate}[label=\alph*)]
        \item Desplazamiento
        \item Trayectoria
        \item Distancia
        \item Vector posición
    \end{enumerate}

    \item \textbf{La aceleración se define fundamentalmente como el cambio en:}
    \begin{enumerate}[label=\alph*)]
        \item La posición por unidad de tiempo
        \item La velocidad por unidad de tiempo
        \item La trayectoria del objeto
        \item La masa del cuerpo
    \end{enumerate}

    \item \textbf{La propiedad de los cuerpos de resistirse a los cambios en su estado de reposo o movimiento se denomina:}
    \begin{enumerate}[label=\alph*)]
        \item Fuerza
        \item Aceleración
        \item Inercia
        \item Velocidad
    \end{enumerate}
\end{enumerate}

\section*{II. SECCIÓN PRÁCTICA (6 Puntos)}

\begin{enumerate}[label=\arabic*.]
    \item \textbf{Un vehículo se desplaza a una velocidad constante de 90 km/h. ¿Cuál es su velocidad expresada en metros por segundo (m/s)?}
    \begin{enumerate}[label=\alph*)]
        \item 15 m/s
        \item 20 m/s
        \item 25 m/s
        \item 30 m/s
    \end{enumerate}

    \item \textbf{Un auto que parte del reposo alcanza una velocidad de 30 m/s en un intervalo de 6.0 segundos. ¿Cuál fue su aceleración media?}
    \begin{enumerate}[label=\alph*)]
        \item 5.0 $m/s^2$
        \item 6.0 $m/s^2$
        \item 180 $m/s^2$
        \item 4.0 $m/s^2$
    \end{enumerate}

    \item \textbf{Se deja caer una piedra desde lo alto de un edificio y tarda 2.0 segundos en tocar el suelo. ¿Cuál es la altura aproximada del edificio (g = 9.8 $m/s^2$)?}
    \begin{enumerate}[label=\alph*)]
        \item 9.8 m
        \item 19.6 m
        \item 39.2 m
        \item 4.9 m
    \end{enumerate}

    \item \textbf{¿Qué fuerza neta se requiere para impartirle una aceleración de 2.0 $m/s^2$ a un bloque de 10.0 kg de masa?}
    \begin{enumerate}[label=\alph*)]
        \item 5.0 N
        \item 12.0 N
        \item 20.0 N
        \item 10.0 N
    \end{enumerate}
\end{enumerate}

\end{document}
